\section{Stage-conditional shift normalization}
\label{sec:method}

\subsection{What the normalization layer leaves in the representation}

Let $z_i\in\bR^d$ be the penultimate representation of epoch $i$, let
$k_i\in\{1,\dots,K\}$ be its sleep stage, and let $r(i)$ be the recording it
belongs to. Write $\cR$ for a recording and $\pi_{\cR,k}$ for the fraction of
its epochs scored as stage $k$. The class-conditional mean of the
representation in a recording is
\begin{equation}
  m_{\cR,k} \;=\; \bE\!\left[z_i \mid r(i)=\cR,\, k_i=k\right],
  \qquad
  m_{\cR} \;=\; \sum_{k}\pi_{\cR,k}\,m_{\cR,k},
  \label{eq:means}
\end{equation}
and the cohort mean is $\mu_{\mathrm{coh}}=\sum_k \pi_k \mu_k$, with $\mu_k$ the
class mean over the training cohort and $\pi_k$ its stage mix.

The decomposition that organises the paper is a definition, not an
assumption. Set
\begin{equation}
  a_{\cR} \;=\; m_{\cR}-\bar\mu_{\cR},
  \qquad
  s_{\cR,k} \;=\; m_{\cR,k}-\mu_k-a_{\cR},
  \qquad
  \bar\mu_{\cR}=\sum_k \pi_{\cR,k}\mu_k .
  \label{eq:decomp}
\end{equation}
Then $m_{\cR,k}=\mu_k+a_{\cR}+s_{\cR,k}$ and, by construction,
$\sum_k \pi_{\cR,k}s_{\cR,k}=0$. The three terms have different meanings.
$\mu_k$ is the stage prototype the classifier was trained on. $a_{\cR}$ is a
recording-wide offset: it is the same vector for every epoch of the night and
is what a batch-normalization layer manufactures when the recording's first
and second moments differ from the cohort's. $s_{\cR,k}$ is the part of the
patient that is not an offset: within a stage, this patient's spindles, slow
waves or alpha differ from the cohort prototype in a way that is stable across
the night. The constraint makes the offset and the conditional signature
identifiable from each other.

\subsection{Marginal adaptation removes the offset and keeps the signature}

The default label-free remedy is to replace the cohort statistics by the
recording's, either by re-estimating every batch-normalization statistic
(AdaBN) or, at the representation, by subtracting the recording's mean shift
$m_{\cR}-\mu_{\mathrm{coh}}$. The next proposition states exactly what that
leaves behind.

\begin{proposition}[Marginal residual]
\label{prop:marginal}
Under the decomposition \eqref{eq:decomp}, marginal correction
$z_i\mapsto z_i-(m_{\cR}-\mu_{\mathrm{coh}})$ gives
\begin{equation}
  \bE\!\left[z_i-(m_{\cR}-\mu_{\mathrm{coh}})\mid r(i)=\cR,k_i=k\right]
  \;=\;\mu_k+s_{\cR,k}+\bigl(\mu_{\mathrm{coh}}-\bar\mu_{\cR}\bigr).
  \label{eq:marginal-residual}
\end{equation}
\end{proposition}
\begin{proof}
Substitute $m_{\cR}=\bar\mu_{\cR}+a_{\cR}$ from \eqref{eq:decomp} into the
conditional mean $m_{\cR,k}=\mu_k+a_{\cR}+s_{\cR,k}$ and cancel $a_{\cR}$.
\end{proof}

Two consequences follow. The conditional signature $s_{\cR,k}$ survives
untouched: marginal adaptation removes a vector that is constant across stages
and cannot remove one that is not. And the correction adds
$\mu_{\mathrm{coh}}-\bar\mu_{\cR}$, which is zero only when the recording's
stage mix equals the cohort's. A recording with an atypical architecture
therefore keeps a residual offset that is a function of \emph{how it slept},
not of who the patient is. Even an oracle marginal correction that knows
$\bar\mu_{\cR}$ and subtracts $m_{\cR}-\bar\mu_{\cR}$ would leave
$\mu_k+s_{\cR,k}$; the conditional signature is out of reach of any operation
that uses one mean per recording. The next corollary states that this is not a
property of a particular estimate of the mean but of the entire family.

\begin{corollary}[No single-vector correction removes the signature]
\label{cor:no-marginal}
Let $g$ be any function of the recording alone, and apply the correction
$z_i\mapsto z_i-g(\cR)$. Then for any two stages $j,k$,
\begin{equation}
\begin{aligned}
  &\bE\!\left[z_i-g(\cR)\mid r(i)=\cR,k_i=j\right] \\
  &\quad-\bE\!\left[z_i-g(\cR)\mid r(i)=\cR,k_i=k\right]
  =(\mu_j-\mu_k)+(s_{\cR,j}-s_{\cR,k}),
\end{aligned}
  \label{eq:invariant}
\end{equation}
independently of $g$.
\end{corollary}
\begin{proof}
The same vector $g(\cR)$ is subtracted from both conditional means, and
$m_{\cR,j}-m_{\cR,k}=(\mu_j-\mu_k)+(s_{\cR,j}-s_{\cR,k})$ by
\eqref{eq:decomp}.
\end{proof}

The corollary is the precise sense in which the signature is out of reach of
the marginal family. Whatever $g$ is --- the recording mean, an oracle that
knows the patient's stage mix, or a vector learned from the recording --- the
displacement the head compares between two stages changes by the same
$s_{\cR,j}-s_{\cR,k}$. A better estimate of the marginal statistic moves the
offset and leaves the relative geometry exactly where it was; only an operation
whose index depends on $k$ can change \eqref{eq:invariant}. The statement is a
consequence of the definition rather than of a modeling assumption: the
conditional means are what they are, and the difference between two of them
contains $s_{\cR,j}-s_{\cR,k}$ exactly. What the corollary rules out is
therefore not a particular estimator but a class of interventions.

The consequence for the classifier is equally direct. The logit margin
between two stages under marginal correction is
\begin{equation}
  W\!\left(m'_{\cR,j}-m'_{\cR,k}\right)
  = W\!\left(\mu_j-\mu_k+s_{\cR,j}-s_{\cR,k}\right),
  \label{eq:marginal-margin}
\end{equation}
so the patient's stage-conditional signature is added to the very quantity
the classifier compares. When that signature is large --- an atypical patient,
a recording with unusual stage-specific morphology --- the margin can move
across the decision boundary, and accuracy falls. This is the mechanism behind
the accuracy difference between the marginal and conditional corrections in
Section~\ref{sec:scsn-identity}.

\subsection{Stage-conditional shift normalization}

The alternative is to estimate the shift separately for each stage and to
remove it with the model's own soft assignments. Let
$p_{ik}=\Pr(k_i=k\mid z_i)$ be the posterior of the unadapted model, and let
\begin{equation}
  \hat m_{\cR,k}
  =\frac{\sum_{i:r(i)=\cR}p_{ik}z_i+\lambda_s\mu_k}
         {\sum_{i:r(i)=\cR}p_{ik}+\lambda_s},
  \qquad
  \hat\delta_{\cR,k}=\hat m_{\cR,k}-\mu_k ,
  \label{eq:estimator}
\end{equation}
where $\lambda_s\ge0$ is a pseudo-count that shrinks rare stages towards the
cohort prototype; in all experiments $\lambda_s=10$. The corrected
representation is
\begin{equation}
  \tilde z_i \;=\; z_i-\sum_{k}p_{ik}\,\hat\delta_{\cR,k}.
  \label{eq:scsn}
\end{equation}
Equations \eqref{eq:estimator}--\eqref{eq:scsn} are stage-conditional shift
normalization (SCSN). The operation is a closed-form post-processing of one
forward pass: the pass gives $p_{ik}$ and $z_i$, the per-stage shifts are
accumulated, and the corrected representation is obtained by subtracting the
weighted shifts. No second pass over the network is required. No weight is
updated, no optimizer is run, no label of the recording is read, and no
subject identifier is used.

Read as a normalization operator, \eqref{eq:scsn} has the same shape as the
normalizers reviewed in Section~\ref{sec:related} but different sufficient
statistics: batch normalization uses the moments of a batch, group
normalization the moments of an epoch, and SCSN the class-conditional means of
the recording, weighted by the model's posterior. It is a normalization whose
statistics are indexed by the classes the model assigns.

For a recording of $n$ epochs and penultimate dimension $d$, accumulating the
stage summaries takes $O(nKd)$ arithmetic and $O(Kd)$ additional storage; the
correction has the same arithmetic complexity. No network replay, backward
pass, or gradient is required.

\begin{proposition}[Conditional residual]
\label{prop:conditional}
Let $\hat\delta_{\cR,k}=\delta_{\cR,k}=a_{\cR}+s_{\cR,k}$ (the population
value). Then for every epoch,
\begin{equation}
  \tilde z_i
  = z_i-\sum_k p_{ik}\bigl(a_{\cR}+s_{\cR,k}\bigr)
  = z_i-a_{\cR}-\sum_k p_{ik}s_{\cR,k} .
  \label{eq:cond-exact}
\end{equation}
In particular the offset cancels for \emph{any} posterior with
$\sum_k p_{ik}=1$, and if $p_{ik}=\mathbb{1}[k_i=k]$ the conditional mean of
$\tilde z_i$ is exactly $\mu_k$.
\end{proposition}
\begin{proof}
$\sum_k p_{ik}=1$ gives $\sum_k p_{ik}a_{\cR}=a_{\cR}$. For hard assignments
$\sum_k p_{ik}s_{\cR,k}=s_{\cR,k_i}$, and $\mu_k+a_{\cR}+s_{\cR,k}-a_{\cR}
-s_{\cR,k}=\mu_k$.
\end{proof}

For a soft posterior the residual is
\begin{equation}
  \bE[\tilde z_i\mid r(i)=\cR,k_i=k]
  =\mu_k+\sum_j\bigl(\mathbb{1}[j=k]-p_{ij}\bigr)s_{\cR,j},
  \label{eq:soft-residual}
\end{equation}
so the error is the part of the conditional signature assigned to the wrong
stage. It vanishes when the posterior is correct and is bounded by
$\max_j\|s_{\cR,j}\|\sum_j|\mathbb{1}[j=k]-p_{ij}|$. This is why the method is
robust to an imperfect but informative posterior and why it degrades when the
posterior is close to uniform, a prediction we test in the cross-cohort
experiment.

\begin{remark}[Why the posterior-weighted average]
\label{rem:mmse}
Condition on the recording and write $\Delta_i=a_{\cR}+s_{\cR,k_i}$ for the
shift that affects epoch $i$. Given $z_i$, the least-mean-square estimate of
$\Delta_i$ among functions of $z_i$ is the conditional expectation
\begin{equation}
  \bE[\Delta_i\mid z_i]=a_{\cR}+\sum_k \Pr(k_i=k\mid z_i)\,s_{\cR,k},
  \label{eq:mmse}
\end{equation}
and \eqref{eq:scsn} subtracts this quantity with the model's posterior in place
of the unknown true one. The correction is therefore not a heuristic blend: it
is the estimate that minimizes the expected squared residual, and the residual
it leaves, \eqref{eq:soft-residual}, is the part of the shift attributed to the
wrong stage. The hard-assignment variant replaces the posterior by a point mass;
it has a smaller conditional bias and a larger variance, and the experiments
show that it over-corrects. The interpolated family
$z_i-\delta_{\cR}-\lambda(\sum_k p_{ik}\hat\delta_{\cR,k}-\delta_{\cR})$
traces a path from ignoring the stage ($\lambda=0$) to trusting the posterior
completely ($\lambda=1$).
\end{remark}

The class geometry is preserved by the same argument. Under the population
correction the class-conditional means of the corrected representation are
$\mu_k$, independently of the recording, so
\begin{equation}
  W\!\left(\tilde m_{\cR,j}-\tilde m_{\cR,k}\right)
  = W\!\left(\mu_j-\mu_k\right)
  \label{eq:scsn-margin}
\end{equation}
for every recording: the logit margin the classifier was trained on is
restored, whereas under marginal correction it is perturbed by
$s_{\cR,j}-s_{\cR,k}$ as in \eqref{eq:marginal-margin}. This is the formal
content of the accuracy result: SCSN does not merely remove a nuisance, it
maps every recording back to the class geometry on which the head was fitted.

Two controls follow from the algebra. If the correction is constrained to the
null space of $W$, the logits are unchanged exactly, so any change in identity
cannot be attributed to a change in predictions; we implement this by
projecting each $\hat\delta_{\cR,k}$ off the row space of $W$. If the shifts of
a \emph{different} recording are applied, the correction is a non-zero
recording-specific vector and the identity is replaced rather than removed;
this is the matched-versus-mismatched control.

\begin{remark}[Estimation error]
With $n_{\cR,k}=\sum_i p_{ik}$ effective epochs, the shrunk estimator
\eqref{eq:estimator} has bias at most
$\lambda_s\|m_{\cR,k}-\mu_k\|/(n_{\cR,k}+\lambda_s)$ towards the cohort
prototype and standard deviation of order $\|s_{\cR,k}\|/\sqrt{n_{\cR,k}}$.
Both terms are small for the stages that dominate a night (N2, REM) and large
for N1, which is why the shrinkage is kept in all experiments. On HMC a
recording averages about 900 epochs, of which roughly 100 are scored N1 and
330 N2, so a pseudo-count of ten moves the N1 estimate by at most a tenth
towards the prototype and the N2 estimate by at most a thirtieth. The estimator
is a soft, regularized version of the stage-conditional sample mean; no
gradient, iteration or held-out label enters it.
\end{remark}

\begin{algorithm}[t]
\caption{Stage-conditional shift normalization (one recording)}
\label{alg:scsn}
\begin{algorithmic}[1]
\Require frozen stager $f=W\phi+b$, unlabeled epochs $x_1,\dots,x_n$ of a recording $\cR$
\Require cohort class means $\mu_k$ and cohort mean $\mu_{\mathrm{coh}}$
\State $z_i\gets\phi(x_i)$ \Comment{forward pass}
\State $p_{ik}\gets\mathrm{softmax}(Wz_i+b)$
\For{$k=1,\dots,K$}
  \State $\hat m_{\cR,k}\gets(\sum_i p_{ik}z_i+\lambda_s\mu_k)/(\sum_i p_{ik}+\lambda_s)$
  \State $\hat\delta_{\cR,k}\gets\hat m_{\cR,k}-\mu_k$
\EndFor
\State $\tilde z_i\gets z_i-\sum_k p_{ik}\hat\delta_{\cR,k}$ \Comment{correction}
\State \Return $\mathrm{softmax}(W\tilde z_i+b)$
\end{algorithmic}
\end{algorithm}

\subsection{Relation to existing normalizers}

\paragraph{Batch normalization and AdaBN}
A batch-normalized stager at inference is an affine map with cohort
statistics. AdaBN replaces those statistics by the recording's at every layer
\citep{li2018adabn}. That is a nonlinear, layer-wise operation and is not the
same object as the representation-level marginal correction of
Proposition~\ref{prop:marginal}; the proposition describes the post-hoc
operation that subtracts one vector from the frozen penultimate
representation, which is the operation SCSN generalises. We report both: the
post-hoc marginal arm isolates the algebra, and layerwise AdaBN is the
stronger inference-time baseline that a practitioner with access to the model
would use. Setting $\lambda=0$ in the interpolated form
$z_i-\delta_{\cR}-\lambda(\sum_k p_{ik}\hat\delta_{\cR,k}-\delta_{\cR})$
recovers the post-hoc marginal correction; $\lambda=1$ is SCSN. We report the
whole $\lambda$ path.

\paragraph{Group and instance normalization}
Group normalization removes the batch from the computation
\citep{wu2018group}: the output of a group-normalized layer depends only on
the epoch it is applied to. Standard AdaBN therefore has no batch statistic to
replace. Standard Tent updates batch-normalization affine parameters; an
entropy objective on the affine parameters of group normalization would be a
different test-time method and is not evaluated here. It does not follow that
the learned representation is free of the recording. Group normalization
removes the per-epoch mean at each layer, but the network that follows maps
the remaining recording-dependent variation into the penultimate space; the
experiments below show that a linear probe still identifies the patient from
a group-normalized representation, and that the post-hoc marginal correction
still changes both accuracy and calibration. The post-hoc marginal correction
can remove only the recording-wide offset, and the experiments show that it
moves the within-stage probe by no more than a third of what SCSN removes.
SCSN is the only stage-aware statistics-replacement method evaluated there,
and we therefore apply it to these checkpoints.

\paragraph{Class-conditional and entropy-based test-time adaptation}
Tent minimizes the entropy of the posterior over the affine parameters of the
normalization layers \citep{wang2021tent}; it is class-agnostic and optimizes a
surrogate. SCSN has a closed form, one hyperparameter (the shrinkage), and a
target: the per-stage mean. The two are compared directly in
Section~\ref{sec:accuracy}.

\paragraph{Subject-invariant architectures}
Adversarial subject classifiers and subject identifiers
\citep{liao2020trifeaturenet,phan2022sleepformer} change the weights and need
retraining; disentangled architectures \citep{perslev2021} do the same. SCSN
is applied to a frozen checkpoint and changes no parameter. It is therefore
not a competitor to those methods but a post-hoc alternative for models that
already exist, and it can be stacked on top of any of them.
